% ===================================================================
%  TAULA N° 13 — L'Ostalariá. — Lo Restaurant. — Lo Cafè.
%  Traduction 2026 d'apres texto/io, controlee sur texto/fr et
%  texto/ca. Consigne : voir texto/oc/10-taula-01.tex.
%
%  Deux notes, marquees toutes deux « (*) », et deux appels : celui du
%  bloc 01-2 (la chambre) et celui du bloc 09-1 (le menu).
% ===================================================================

%%K t13-noto-1 noto
\VUnotes{143.2mm}{%
\textit{\VUgras{(*) Pels detalhs de la cambra, vejatz la taula n° 6.}}%
}

%%K t13-apar-1 apar
\VUsaut{3.0mm}
\VUtitre{15.0pt}{60}{TAULA N° 13}
\VUsaut{1.6mm}
\VUpk{10.2pt}{40}{L'Ostalariá. --- Lo Restaurant.}
\VUsaut{2.0mm}
\VUpk{10.2pt}{40}{Lo Cafè.}
\VUsaut{3.4mm}

%%K t13-01-1 p
1. --- Essent arribat a Royan ièr al vèspre, m'alotgèri a l'\textit{Ostalariá de
l'Ocean}, que me l'aviá recomandada un amic.

%%K t13-01-2 p
Me donèron una bèla \VUgras{cambra} \textsuperscript{(2)} al segond
\VUgras{estatge} \textsuperscript{(1)}, ont dormiguèri fòrça plan (*).

%%K t13-02-1 p
2. --- Aqueste matin, en sortir de ma cambra, ont la \VUgras{cambrièira}
\textsuperscript{(3)} m'aviá portat lo dejunar, dintrèri dins lo \VUgras{salon de
fumar} \textsuperscript{(5)}, que sèrv tanben de \VUgras{sala de lectura}
\textsuperscript{(4)}, per fulhejar los \VUgras{jornals} \textsuperscript{(6)}
mentre que fumavi una \VUgras{cigareta} \textsuperscript{(8)}. Quand aguèri
acabat de legir, davalèri amb l'\VUgras{ascensor} \textsuperscript{(9)} e sortiguèri
far un torn per la vila.

%%K t13-03-1 p
3. --- Coma m'èra escapat de demandar a l'\VUgras{emplegat}
\textsuperscript{(12)} de l'ostalariá a quina ora se servissián los repais a la
\VUgras{taula redonda} \textsuperscript{(11)}, tornèri d'ora e o profitèri per me
donar un banh dins la \VUgras{sala de banh} \textsuperscript{(10)} de
l'ostalariá. Puèi anèri a la \VUgras{caissa} \textsuperscript{(15)} per telefonar
a un de mos amics. Mas lo \VUgras{telefòn} \textsuperscript{(13)} èra ocupat; en
veire mon embarràs, la \VUgras{caissièira} \textsuperscript{(14)}, qu'èra a marcar
un \VUgras{compte} \textsuperscript{(17)} sul \VUgras{registre dels viatjaires}
\textsuperscript{(16)}, demandèt al \VUgras{concierge} \textsuperscript{(18)}
que mandèsse lo \VUgras{grom} \textsuperscript{(19)} far ma comission amb sa
\VUgras{bicicleta} \textsuperscript{(20)}.

%%K t13-04-1 p
4. --- Li mercegèri e sortiguèri donar una estrena al \VUgras{portafais}
\textsuperscript{(24)} (l'òme per tot), qu'aviá portat de l'estacion sus sa
\VUgras{bèrla} \textsuperscript{(25)} ma \VUgras{bóstia de capèls}
\textsuperscript{(26)}, ma \VUgras{valisa} \textsuperscript{(27)} e ma
\VUgras{saca de viatge} \textsuperscript{(28)}. Lo \VUgras{factor}
\textsuperscript{(21)}, qu'arribava en aquel moment, me liurèt una \VUgras{letra}
\textsuperscript{(22)}.

%%K t13-05-1 p
5. --- Demorèri una estona de mai sul lindal, contra la \VUgras{bóstia de las
letras} \textsuperscript{(23)}, a agachar lo movement del carrièr. Lo
\VUgras{conductor} \textsuperscript{(30)} de l'\VUgras{omnibús}
\textsuperscript{(29)} cargava sus son vehicul la \VUgras{mala}
\textsuperscript{(31)} d'un \VUgras{viatjaire} \textsuperscript{(32)} al qual
l'\VUgras{ostalièr} \textsuperscript{(33)} disiá adieu. Un jove
\VUgras{telegrafista} \textsuperscript{(35)} portava, sens cap de pressa, un
\VUgras{telegrama} \textsuperscript{(34)} per un dels òstes de l'ostalariá; un
\VUgras{torista} \textsuperscript{(36)}, amb sa \VUgras{camèra fotografica}
\textsuperscript{(38)} pendolada a l'espatla, consultava sa \VUgras{guida}
\textsuperscript{(37)}.

%%K t13-tit-1 sub
\VUsaut{4.0mm}
\VUpk{10.2pt}{40}{L'ora del dinnar.}
\VUsaut{3.0mm}

%%K t13-06-1 p
6. --- Davant la reissa, un plàcid \VUgras{comissionari} \textsuperscript{(39)}
penequejava entre sos \VUgras{cròcs de carga} \textsuperscript{(40)} e sa
\VUgras{bóstia de lustraire} \textsuperscript{(41)}, mentre qu'una jove
\VUgras{bugadièira} \textsuperscript{(42)} que portava un \VUgras{panièr}
\textsuperscript{(43)} de linja risiá de l'aire solemne del \VUgras{cap de
servici} \textsuperscript{(44)} plantat contra la pòrta.

%%K t13-07-1 p
7. --- La vista del \VUgras{cosinièr} \textsuperscript{(45)}, qu'èra sortit de sa
\VUgras{cosina} \textsuperscript{(47)} del \VUgras{sosterranh}
\textsuperscript{(48)} per mandar un \VUgras{marmiton} \textsuperscript{(46)} a
una comission, me remembrèt que s'apropchava l'ora de dinnar, e anèri al restaurant
de l'autre costat de la cort, ont un \VUgras{automobilista}
\textsuperscript{(50)} garava son \VUgras{automobila} \textsuperscript{(49)}
mentre qu'un \VUgras{mecanician} \textsuperscript{(52)} adobava, segon las
indicacions del \VUgras{ciclista} \textsuperscript{(53)}, un \VUgras{tricicle de
benzina} \textsuperscript{(51)} que veniá d'aver un accident.

%%K t13-08-1 p
8. --- Demandèri a un jove \VUgras{grom} \textsuperscript{(54)} que portava de
\VUgras{gèrlas de cervesa} \textsuperscript{(55)} se la taula redonda èra jos la
veranda, e coma me diguèt que òc, m'assetèri a una de las taulas e me serviguèron
un repais excellent.

%%K t13-noto-2 noto
\VUnotes{143.2mm}{%
\textit{\VUgras{(*) Amb l'ajuda de la lista de plats que dona lo vocabulari, lo
mèstre pòt variar aisidament lo menut seguent.}}%
}

%%K t13-tit-2 sub
\VUsaut{4.0mm}
\VUpk{10.2pt}{40}{Un dinnar excellent.}
\VUsaut{3.0mm}

%%K t13-09-1 p
9. --- Lo garçon me portèt lo menut del dinnar (*) e la carta dels vins.
Causiguèri e demandèri de \VUgras{gambas} amb de \VUgras{burre} coma
\VUgras{entremés}, e, coma primièr plat, de \VUgras{tripas a la mòda de Caen}.
Prenguèri pas de \VUgras{peis}, mas demandèri una racion de \VUgras{cuèissa} de
moton e, coma \VUgras{legum}, d'\VUgras{espinarcs} a la crèma. Puèi, coma cap dels
\VUgras{plats doces} me tentava pas, faguèri portar un assortiment de postres:
de \VUgras{fromatge}, de \VUgras{frucha} e de \VUgras{gaufretas}. I apondèri un
excellent \VUgras{vin} de Sant Emilion e, fòrça satisfach de mon repais, me levèri
de taula aprèp aver donat al \VUgras{garçon} \textsuperscript{(57)} una bona
\VUgras{estrena} \textsuperscript{(59)}.

%%K t13-10-1 p
10. --- En veire que m'aprestavi a partir, lo \VUgras{gerent}
\textsuperscript{(60)} me prepausèt de sortir sul balcon per admirar l'esplendid
panorama de l'\VUgras{ocean} \textsuperscript{(62)}. Al luènh se destriavan de
pichonas \VUgras{barcas} \textsuperscript{(63)} de pesca, de \VUgras{velièrs}
\textsuperscript{(64)} e un enòrme \VUgras{transatlantic} \textsuperscript{(65)}
que sa silueta negra se retalhava suls \VUgras{nívols} \textsuperscript{(67)}
blancs de l'\VUgras{orizont} \textsuperscript{(66)}. Una brisa leugièra agitava la
\VUgras{bandièra} \textsuperscript{(70)} plantada dessús l'\VUgras{ensenha}
\textsuperscript{(69)} del \VUgras{vestibul} \textsuperscript{(68)}.

%%K t13-tit-3 sub
\VUsaut{3.0mm}
\VUpk{10.2pt}{40}{Divertiments.}
\VUsaut{3.0mm}

%%K t13-11-1 p
11. --- L'espectacle èra tan interessant que decidiguèri de pujar l'endeman a la
terrassa del cafè, en portar amb ieu de \VUgras{binòcles} \textsuperscript{(71)} o
una \VUgras{lunèta} \textsuperscript{(72)}.

%%K t13-12-1 p
12. --- En sortir del balcon anèri al cafè de l'ostalariá prene una
\VUgras{consomacion} \textsuperscript{(73)} e jogar una partida de
\VUgras{bilhard} \textsuperscript{(75)}. Traversèri la sala del cafè, ont fòrça
clients jogavan als \VUgras{escacs} \textsuperscript{(78)}, a las \VUgras{damas}
\textsuperscript{(79)}, al \VUgras{dominò} \textsuperscript{(80)}, al
\VUgras{tricatrac} \textsuperscript{(81)} o a las \VUgras{cartas}
\textsuperscript{(82)}, e pugèri drech a la \VUgras{sala de bilhard}
\textsuperscript{(74)}.

%%K t13-13-1 p
13. --- Acabada la partida, davalèri al ras-de-caussada e demandèri al garçon una
\VUgras{envelòpa} \textsuperscript{(84)}, de \VUgras{papièr}
\textsuperscript{(83)}, un \VUgras{timbre} \textsuperscript{(85)}, una
\VUgras{pluma} \textsuperscript{(87)} e de \VUgras{tinta} \textsuperscript{(86)}
per escriure una letra.

%%K t13-13-2 p
Puèi sonèri un \VUgras{cochièr} \textsuperscript{(92)} que sa \VUgras{veitura}
\textsuperscript{(88)} s'èra arrestada davant lo cafè. Li demandèri son prètz per
ora e per corsa e, un còp d'acòrdi, me dobriguèt la \VUgras{portièra}
\textsuperscript{(89)} e m'installèt dedins. Tornèt pujar sus son \VUgras{sèti}
\textsuperscript{(91)}, butèt vivament los cavals amb son \VUgras{foet}
\textsuperscript{(93)}, e cap enlà anèrem far un passejament deliciós.
\VUsaut{6.0mm}
\VUfilet{22mm}
